\chapter{Defect lines, twisted sectors, and boundaries}
\label{ch:defects}

A topological line can be wrapped around an operator, inserted in
a thermal trace, or drawn along the time direction to define a
twisted Hilbert space. A conformal boundary can be viewed either
as a state in a closed-channel Hilbert space or as an endpoint
condition for an open-channel Hamiltonian. Modular transformations
relate these descriptions. The Haagerup examples make all of these
relations explicit and show why six lines, twelve central sectors,
and thirty-six bimodule defects can coexist without contradiction.

\section{The eigenvalue of a contractible line}
\index{defect!Verlinde line}\index{invariant projector}

For an ordinary symmetry element, acting on the vacuum gives one.
For a general simple topological line $a$, a contractible vacuum
loop evaluates to its quantum dimension $d_a$. Therefore the
appropriate invariance condition on a local state is
\begin{equation}
 \widehat L_a|\psi\rangle=d_a|\psi\rangle
 \quad\text{for every symmetry line }a.
 \label{eq:defect-invariant}
\end{equation}
Demanding eigenvalue one for a noninvertible line would generally
discard the vacuum itself.

In a rational diagonal theory, the chiral-algebra-preserving
Verlinde line labelled by $a$ acts on the bulk sector
$M_i\otimes\overline{M_i}$ by
\begin{equation}
 \lambda_a(i)=\frac{S_{ai}}{S_{0i}}.
 \label{eq:defect-eigenvalue}
\end{equation}
At $i=0$ this is $d_a$. Fusion of line operators follows from
the Verlinde formula because the columns of $S$ diagonalize all
fusion matrices simultaneously.

For the Haagerup center all labels are self-dual and $S$ is real.
The Verlinde coefficients are
\begin{equation}
 N_{ab}^{\ c}=\sum_i\frac{S_{ai}S_{bi}S_{ci}}{S_{0i}}.
 \label{eq:defect-Verlinde}
\end{equation}
They are nonnegative integers, even though individual line
eigenvalues need not be positive. The value of a trace with a
line inserted is therefore not a state degeneracy.

\begin{example}[An invariant projector from modular orthogonality]
Let $\mathcal D=1/S_{00}=3(1+d^2)$. On the diagonal bulk space,
\[
 \Pi_{\rm inv}=\frac1{\mathcal D^2}
                       \sum_a d_a\widehat L_a.
\]
Its eigenvalue on sector $i$ is
\[
 \frac1{\mathcal D^2S_{00}S_{0i}}
          \sum_a S_{a0}S_{ai}
 =\frac{\delta_{0i}}{\mathcal D^2S_{00}S_{0i}}
 =\delta_{0i}.
\]
Thus the invariant space is exactly $U\otimes\overline U$.
Equivalently, if all line eigenvalues in sector $i$ equalled
their vacuum values, column $i$ of $S$ would be proportional to
the vacuum column, forcing $i=0$.
\end{example}

\section{Twining and twisting are different traces}
\index{twined trace}\index{twisted sector}

Insert the line operator in the ordinary Hilbert-space trace:
\begin{equation}
 Z^{[a]}(\tau,\bar\tau)=
 \sum_i\lambda_a(i)|\chi_i(\tau)|^2.
 \label{eq:defect-twined}
\end{equation}
This is a \emph{twined} partition function. Its coefficients can
be negative or algebraic because they are traces of a symmetry
operator on states.
Exchange the temporal and spatial cycles of the torus. The line
now intersects the spatial circle, and the corresponding
\emph{twisted} Hilbert space has partition function
\begin{equation}
 Z_{[a]}(\tau,\bar\tau)=
 \sum_{i,j}N_{ai}^{\ j}\chi_i(\tau)\overline{\chi_j(\tau)}.
 \label{eq:defect-twisted}
\end{equation}
To prove this relation, transform each character in
\eqref{eq:defect-twined} by $S$. The coefficient of
$\chi_i\overline{\chi_j}$ is
$\sum_kS_{ak}S_{ik}S_{jk}/S_{0k}$, exactly the Verlinde sum.
Here every coefficient is a nonnegative integer because a
twisted Hilbert space counts states.

The ground dimension in a twisted sector need not equal its
first scalar dimension. A term
$M_i\otimes\overline{M_j}$ has scalar states only if
$h_i-h_j\in\Z$ and the grades can compensate that difference.
For the present twelve lowest weights, the only unequal values
differing by an integer are zero and one. Thus the lowest scalar
count can be obtained from the top dimensions and the six
vacuum currents without expanding many character grades.

\begin{example}[A light twisted ground state need not be a scalar]
The line $c_1$ has twisted ground dimension $1/3$, with six
components from the left/right placements of its three top states
against the vacuum. Their spins are $\pm1/3$.
Its first scalar dimension is instead $2/3$, with multiplicity
$2\cdot3^2=18$. We can derive the factor two from the
modular matrix instead of reading a fusion table. Put
$A=S_{00}$, $B=S_{0b}$, so that $AB=1/117$ and
$A+B=1/3$. The nonzero entries in the $c_1$ row are
the first two $1/3$ entries and a middle block consisting
of one $2/3$ and three $-1/3$ entries. Verlinde gives
\[
 N_{c_1c_1}^{\ c_1}
 =\frac1{27}\left(\frac1A+\frac1B\right)
   +3\left[\left(\frac23\right)^3
                +3\left(-\frac13\right)^3\right]
 =\frac{13}{9}+\frac59=2.
\]
Pairing the three left top states with the three right top
states then gives eighteen. Fractional spin is permitted for states attached
to the nontrivial defect; it would not describe an ordinary
bosonic local field in the untwisted theory.
\end{example}

For orientation, some exact twisted gaps are
\begin{center}
\begin{tabular}{c|cc|cc}
\toprule
Line&Ground $\Delta$&Multiplicity&First scalar $\Delta$&Multiplicity\\
\midrule
$0$&0&1&$2/3$&9\\
$b$&$2/3$&18&$2/3$&18\\
$a,c_0$&$2/3$&9&$2/3$&9\\
$c_1$&$1/3$&6&$2/3$&18\\
$c_2$&$2/3$&63&$2/3$&9\\
$d_1$&$6/13$&14&$2/3$&9\\
$d_2,d_3$&$2/3$&9&$2/3$&9\\
$d_4$&$5/13$&10&$2/3$&9\\
$d_5$&$7/13$&26&$2/3$&9\\
$d_6$&$8/13$&28&$2/3$&9\\
\bottomrule
\end{tabular}
\end{center}
The vacuum is excluded from the phrase ``first scalar'' in the
identity row. The common scalar gap does not make the twisted
spaces equivalent; their ground states, spins, and multiplicities
already distinguish them.

\section{Cardy boundaries and open--closed duality}
\index{Cardy boundary}\index{open--closed duality}\index{boundary entropy}

A conformal boundary preserving $U$ glues the holomorphic and
antiholomorphic symmetry. For the stress tensor, the boundary
state obeys
\[
 (L_n-\bar L_{-n})|a\rangle=0.
\]
The opposite mode number comes from the reversal of orientation
at the boundary. For each compatible sector there is an
Ishibashi state $|i\rangle\!\rangle$, a formal solution of the
gluing equations. It is not a normalizable vector before a
positive Euclidean propagation length is inserted. It pairs
matching left and right basis states in the module. With the
standard normalization, the regularized inner product is
\[
 \langle\!\langle i|
 \tilde q^{(L_0+\bar L_0-c/12)/2}
 |j\rangle\!\rangle=\delta_{ij}\chi_i(\tilde q).
\]
Each matching basis pair contributes one power
$\tilde q^{h_i+n-c/24}$, which explains the character.

The annulus admits two quantizations
(Figure~\ref{fig:defect-annulus-channels}).
On an interval of length $\ell$, with Euclidean time period
$\beta$, the open-channel Hamiltonian is
$H_{\rm open}=(\pi v/\ell)(L_0-c/24)$ and
$q=e^{-\pi v\beta/\ell}$. Quantizing across the interval
instead gives propagation between boundary states on a circle,
with $\tilde q=e^{-4\pi\ell/(v\beta)}$. These two nomes are
related by the modular $S$ transformation. A consistent boundary
condition must make the closed-channel overlap into a trace with
nonnegative integer open-channel multiplicities.

The factors of two follow from the geometry. Unfolding the open
interval gives one chiral circle of length $2\ell$, hence the
coefficient $\pi v/\ell$ in $H_{\rm open}$. In the closed channel,
the spatial circumference is $v\beta$. Propagation over distance
$\ell$ therefore inserts
\[
 \exp\!\left[-\frac{2\pi\ell}{v\beta}
                   (L_0+\bar L_0-c/12)\right],
\]
which is exactly the Ishibashi regulator above with the stated
$\tilde q$. This also explains why the open trace is a sum of chiral
characters rather than absolute squares of characters.

\begin{figure}[htbp]
\centering
\begin{tikzpicture}[>=Stealth,font=\small]
\foreach \x in {0,6}{
 \draw[very thick,MidnightBlue] (\x,0)--(\x,2);
 \draw[very thick,MidnightBlue] (\x+3.5,0)--(\x+3.5,2);
 \draw[dashed] (\x,0)--(\x+3.5,0);
 \draw[dashed] (\x,2)--(\x+3.5,2);
 \node[above] at (\x+1.75,2) {periodic Euclidean time};
 \node[left] at (\x,1) {$a$};
 \node[right] at (\x+3.5,1) {$b$};
}
\draw[->,thick] (1.75,.3)--(1.75,1.7);
\node[below,align=center] at (1.75,-.2)
 {Open channel\\trace over $\mathcal H_{ab}$};
\draw[->,thick] (6.5,1)--(9,1);
\node[below,align=center] at (7.75,-.2)
 {Closed channel\\propagation between boundary states};
\end{tikzpicture}

\caption{The same annulus viewed in two directions. The dashed
horizontal edges are identified. The solid vertical edges carry
boundary conditions $a,b$. The open channel propagates states
along periodic time; the closed channel propagates from one
boundary state to the other.}
\label{fig:defect-annulus-channels}
\end{figure}

The diagonal
theory's elementary Cardy boundaries are
\begin{equation}
 |a\rangle=\sum_i\frac{S_{ai}}{\sqrt{S_{0i}}}
                                      |i\rangle\!\rangle.
 \label{eq:defect-Cardy}
\end{equation}
Their closed-channel overlap contains coefficients
$S_{ai}S_{bi}/S_{0i}$. Applying a modular transformation gives
the open-channel annulus
\begin{equation}
 Z_{ab}^{\rm open}(q)=\sum_iN_{ab}^{\ i}\chi_i(q).
 \label{eq:defect-annulus}
\end{equation}
The usual conjugate boundary label has disappeared only because
the twelve center labels are self-dual. It will be necessary
again for the six Haagerup boundary labels.

The vacuum overlap is the boundary $g$ factor:
\begin{equation}
 g_a=\frac{S_{a0}}{\sqrt{S_{00}}}
     =\frac{d_a}{\sqrt{\mathcal D}},
 \qquad \log g_a\text{ is the boundary entropy}.
 \label{eq:defect-g}
\end{equation}
It is not an integer ground-state degeneracy. In the Haagerup
diagonal theory its four possible values occur in the ratios
$1:d^2:s:3d$.

\begin{example}[Fixed and free Ising boundaries]
\label{ex:defect-Ising-boundaries}
For the Ising labels $(\one,\epsilon,\sigma)$,
\[
 S=\frac12
 \begin{pmatrix}
 1&1&\sqrt2\\
 1&1&-\sqrt2\\
 \sqrt2&-\sqrt2&0
 \end{pmatrix}.
\]
The fusion rules are
$\epsilon^2=\one$, $\epsilon\sigma=\sigma$,
and $\sigma^2=\one+\epsilon$. The boundaries labelled
$\one,\epsilon$ are the two fixed-spin boundaries, and
$\sigma$ is the free boundary. Cardy's formula gives
\[
 Z_{++}=\chi_{\one},\qquad
 Z_{+-}=\chi_\epsilon,\qquad
 Z_{ff}=\chi_{\one}+\chi_\epsilon,\qquad
 Z_{+f}=\chi_\sigma.
\]
For instance, the free boundary has one nonidentity relevant boundary
primary of weight $1/2$, while a boundary-changing field
between a fixed and a free boundary has weight $1/16$.
The vacuum row gives
$g_+=g_-=1/\sqrt2$ and $g_f=1$. The noninteger value
$1/\sqrt2$ is entirely consistent because $g$ is a vacuum
overlap, not an open-channel ground multiplicity.
\end{example}

\begin{example}[The vacuum and $b$ boundary annuli]
The vacuum boundary has $Z_{00}^{\rm open}=\chi_0$. Its first
positive boundary weight is one, with six currents.
For $b$, the fusion rule
\[
 b^2=\one+b+a+c_0+c_1+c_2+\sum_{j=1}^6d_j
\]
puts three weight-$1/3$ states in $Z_{bb}^{\rm open}$, coming
from the top of $c_1$. They are relevant boundary perturbations.
One must not remove the entire vacuum representation when asking
for a first positive boundary energy: doing so would discard the
currents at the vacuum boundary.
\end{example}

\section{A nine-object example of condensation}
\index{condensation}\index{local module}\index{bimodule}

Before applying the extension dictionary to Haagerup symmetry,
consider the pointed center $Z(\Vec_{\Z_3})$ from
Chapter~\ref{ch:centers}. Its simples are $(g,k)$ with
$g,k\in\Z_3$, fusion adds both labels, and the double-braiding
phase of $(g,k)$ with $(h,\ell)$ is
$\omega^{kh+\ell g}$, where $\omega=e^{2\pi i/3}$.
Take the algebra
\[
 A=(0,0)\oplus(0,1)\oplus(0,2),
\]
whose multiplication adds the second label. This algebra
condenses the three pure charges. The construction is an
ordinary finite Fourier problem because every component is
one-dimensional and the relevant associator is trivial.

For each $g$, define the right module
\[
 x_g=(g,0)\oplus(g,1)\oplus(g,2).
\]
The right action of $(0,\ell)$ sends the $k$ component
to the $k+\ell$ component. These three modules are simple
as $A$-modules: a submodule containing any one component
must contain all three. Every simple right module has one
of these fixed fluxes, so the list is complete.
The balanced tensor product identifies the three equivalent
charge placements and gives
\[
 x_g\otimes_Ax_h=x_{g+h}.
\]
Thus the category of all modules is $\Vec_{\Z_3}$.

Locality imposes an extra requirement on the action. The
double braiding with every condensed component $(0,\ell)$
has phase $\omega^{\ell g}$. Since the action is nonzero
on all components, it is unchanged by this phase only when
$g=0$. Hence the local module category has just its unit
$x_0=A$.

There are more bimodules. One can twist the left action by
any character $\ell\mapsto\omega^{t\ell}$ of $\Z_3$,
independently for each $x_g$. The two labels $(g,t)$ give
nine simple bimodules. This small calculation displays the
three different answers to three different questions:
one local simple, three all-module simples, and nine
bimodule simples. A holomorphic extension can therefore
have a trivial local representation category and still
carry nontrivial boundary and defect categories preserving
the smaller chiral algebra.

\section{Where the six Haagerup lines appear in the extension}
\index{defect!Haagerup action}\index{tube multiplicity}

The twelve Verlinde lines just discussed act on diagonal
sectors $M_i\otimes\overline{M_i}$. In the extension construction
below, the same twelve central labels index tube multiplicity
spaces with dynamical sectors $M_i\otimes\overline{H'}$.
The shared categorical labels do not identify these Hilbert
spaces or their partition functions.

Let $A$ be the selected canonical extension algebra and
$\mathcal C=\Rep(U)_A\simeq\mathcal H_3$ the category of all
right $A$-modules. Local $A$-modules form the trivial category
because the extension is holomorphic. The nonlocal modules remain
essential: discarding them would discard the six Haagerup objects.

The full $E_8$ theory has the corresponding $U$-preserving
bimodule defects
$\mathcal C\boxtimes\mathcal C^{\rm op}$, with $36$ simple
ordered pairs. This counts the rational defects preserving
the specified $U$ structure; it is not a classification of all
defects preserving only the stress tensor.
Here $\mathcal C^{\rm op}$ reverses tensor order.
The displayed equivalence is the canonical-algebra bimodule theorem.
The following calculation illustrates why one factor reverses:
let
$T_{x,y}(m)=x\otimes m\otimes y$. Then
\[
 T_{x,y}T_{x',y'}(m)
 =(x\otimes x')\otimes m\otimes(y'\otimes y).
\]
Thus ordinary composition keeps the left order and reverses the
right order. There is a further convention when comparing directly
with right $A$-modules: a bimodule $B$ acts by
$T_B(M)=M\otimes_A B$, and
\[
 T_B\circ T_C\simeq T_{C\otimes_A B}.
\]
This correspondence reverses composition. Passing to the opposite
endofunctor category and exchanging the two factors recovers the
displayed $\mathcal C\boxtimes\mathcal C^{\rm op}$ bimodule
category. The $T_{x,y}$ calculation alone should therefore not be
read as an identification with unchanged tensor conventions.
Choosing the first factor gives six lines acting
on the holomorphic side and leaving the antiholomorphic $H'$
factor untouched. These defects preserve the original stress,
because $U$ and $H'$ have the same conformal vector. They need
not preserve every enlarged $E_8$ current. This is compatible
with the trivial local module category of the maximal
holomorphic algebra.

Let $F:Z(\mathcal C)\to\mathcal C$ forget the half-braiding,
and abbreviate $X=\rho\oplus\alpha\rho\oplus\alpha^2\rho$.
The forgetful multiplicities are
\begin{align}
 F(0)&=\one,&F(b)&=\one\oplus X,&F(a)&=2\one\oplus X,\notag\\
 F(c_j)&=\alpha\oplus\alpha^2\oplus X,&
 F(d_j)&=X.
 \label{eq:defect-forgetful}
\end{align}
Write $n_{x\mu}=\dim\Hom_{\mathcal C}(F(M_\mu),x)$.
Free--forgetful adjunction proves that restriction of boundary
object $x$ to $U$ is
\[
 Rx\cong\bigoplus_\mu n_{x\mu}M_\mu.
\]
Under the rational defect construction for this selected
extension, the first-factor twisted Hilbert space is
\begin{equation}
 \HH_x\cong\bigoplus_\mu W_x^\mu\otimes\mathcal V_\mu,
 \qquad W_x^\mu=\Hom(F(M_\mu),x),
 \quad\mathcal V_\mu=M_\mu\otimes\overline{H'}.
 \label{eq:defect-tube-decomposition}
\end{equation}
The finite tube factor $W_x^\mu$ and the infinite-dimensional
CFT multiplicity space $\mathcal V_\mu$ have different roles.
For instance $\dim W_1^a=2$; the corresponding untwisted
projector image contains two copies of $\mathcal V_a$.
At left grade $k$ and right grade $\ell$ in a specified $\mu$
summand, the number of states is
$n_{x\mu}a_\mu(k)a_{E_8}(\ell)$. Thus $n_{1a}=2$ and
$a_a(0)=81$ give $162$ holomorphic weight-one states against
the antiholomorphic vacuum. The small tube matrix acts on the
two-copy factor, not on a two-state physical Hilbert space.

The twelve sector functions and the six twisted traces are
\begin{equation}
 Z_\mu^{\rm tube}=\chi_\mu\overline{\chi_{E_8}},
 \qquad Z_x=\sum_\mu n_{x\mu}Z_\mu^{\rm tube}.
 \label{eq:defect-sector-functions}
\end{equation}
The left and right central-charge phases cancel, so the twelve
functions transform by $(S,\Theta)$ rather than
$(S,T_{\rm ch})$. For $x=1$, the formula reduces to the ordinary
$E_8$ partition function.

\section{An explicit noninvertible action on untwisted states}

On the multiplicity spaces of
$H'=U\oplus b\oplus(\C^2\otimes a)$, choose bases so that
\begin{equation}
\begin{array}{c|cc}
 &\widehat\alpha&\widehat\rho\\\hline
 U&1&d\\
 b&1&-d^{-1}\\
 \C^2\otimes a&
 \begin{pmatrix}\omega&0\\0&\omega^2\end{pmatrix}&
 \begin{pmatrix}0&1\\1&0\end{pmatrix}
\end{array}
\qquad\omega=e^{2\pi i/3}.
 \label{eq:defect-H3-action}
\end{equation}
The last matrices satisfy
$\alpha\rho=\rho\alpha^{-1}$. Their three noninvertible
matrices sum to zero, so $\rho^2=1$ on that block agrees with
$\rho^2=1+\rho+\alpha\rho+\alpha^2\rho$.
On the neutral blocks the same rule is the quadratic equation
with roots $d,-1/d$.

Taking traces gives
\begin{align}
 Z^{[\alpha]}&=(\chi_0+\chi_b-\chi_a)
                                     \overline{\chi_{E_8}},\notag\\
 Z^{[\rho]}&=(d\chi_0-d^{-1}\chi_b)
                                     \overline{\chi_{E_8}}.
 \label{eq:defect-E8-twined}
\end{align}
The absence of $\chi_a$ in the second line is the zero trace of
the off-diagonal $2\times2$ matrix. Squaring the holomorphic
factors in absolute value would describe a different paired
left/right insertion, not this selected first-factor action.

\begin{example}[Recovering twisted spectra from the line action]
Write the holomorphic coefficient vectors of the two twined
traces in the center ordering as
\[
 t_\alpha=(1,1,-1,0,\ldots,0),\qquad
 t_\rho=(d,-d^{-1},0,\ldots,0).
\]
Multiplication by the explicit $S$ matrix gives
\[
 St_\alpha=(0,0,0,1,1,1,0,\ldots,0),\qquad
 St_\rho=(0,1,1,1,1,1,1,\ldots,1).
\]
For the second vector, the vacuum entry vanishes because
$S_{0b}=d^2S_{00}$. The middle entries are
$(d-d^{-1})/3=1$, and the six final entries are
$(d+d^{-1})/\sqrt{13}=1$.
The $b$ entry is one as well, using
$d^2-d^{-2}=3\sqrt{13}$.
These are exactly the forgetful multiplicities of $\alpha$
and $\rho$ in \eqref{eq:defect-forgetful}. Thus the small
line matrices, modular transformation, and actual restriction
give three compatible descriptions of the same twisted traces.
\end{example}

Every antiholomorphic $E_8$ weight is integral. Hence the
$c_1,c_2,d_j$ tube multiplicity spaces contain no scalar states
at any grade: their left weights are nonintegral modulo one.
The first invariant scalars are instead in
$U_1\otimes\overline{H'_1}$, with dimension two and multiplicity
$6\cdot248=1488$. This derives the invariant gap quoted earlier.

\section{Six boundaries of the extended theory}
\index{annulus partition function}\index{boundary spectrum}

The same module category gives six elementary boundaries preserving
$U$. For boundary labels $x,y\in\mathcal C$,
\begin{equation}
 Z_{xy}^{E_8,\rm open}
 =\sum_{z,\mu}N_{x^*y}^{\ z}\,n_{z\mu}\chi_\mu.
 \label{eq:defect-E8-annulus}
\end{equation}
Now the dual is essential: $\alpha^*=\alpha^2$, whereas the
three noninvertible simples are self-dual.
At the level of morphisms, duality and the central half-braiding
give
\[
 \Hom_{\mathcal C}(F(M_\mu)\otimes x,y)
 \cong\Hom_{\mathcal C}(F(M_\mu),x^*\otimes y).
\]
Free--forgetful adjunction then identifies this space with the
$M_\mu$ multiplicity in the restriction of $x^*\otimes_Ay$.
This derives the dual in the annulus formula from the boundary
junction rather than guessing it from a commutative fusion table.
The identity multiplicity is consequently
$\dim\Hom(x,y)=\delta_{xy}$ for elementary boundaries.
For example, the $\alpha$ self-annulus uses
$\alpha^*\alpha=\one$; omitting the dual would incorrectly
use $\alpha^2$ and remove its boundary identity.
Restriction multiplies quantum dimensions by
$\dim A=\mathcal D$. Taking the closed-channel vacuum
coefficient gives
$S_{00}\dim R(x^*\otimes_Ay)=d_xd_y$.
Reflection positivity selects positive overlaps, so
\[
 (g_1,g_\alpha,g_{\alpha^2},g_\rho,g_{\alpha\rho},g_{\alpha^2\rho})
 =(1,1,1,d,d,d).
\]

\begin{example}[A noninvertible boundary self-spectrum]
For $x=\rho$, use $\rho^2=1+\rho+\alpha\rho+\alpha^2\rho$.
The three noninvertible objects have identical restriction
characters. Thus
\[
 Z_{\rho\rho}^{E_8,\rm open}
 =\chi_0+4\chi_b+5\chi_a
        +3(\chi_{c_0}+\chi_{c_1}+\chi_{c_2})
        +3\sum_{j=1}^6\chi_{d_j}.
\]
The identity occurs once and the first positive weight is $1/3$,
with $3\cdot3=9$ states. By contrast, the invertible boundary
self-annuli equal $\chi_{E_8}$ and start at positive weight one
with $248$ currents. The different boundary entropies and relevant
boundary spectra make the categorical distinction observable.
\end{example}

\section{Exercises}

\begin{exercise}\label{ex:defect-projector}
Check directly that the invariant projector defined above is
idempotent and Hermitian in the diagonal theory. Why is the
denominator $\mathcal D^2$ rather than $\mathcal D$?

\par\smallskip\noindent{\footnotesize\normalfont
\hyperref[sol:ex:defect-projector]{Solution on page~\pageref*{sol:ex:defect-projector}.}}\par
\end{exercise}

\begin{exercise}\label{ex:defect-twined}
Using the matrices in \eqref{eq:defect-H3-action}, verify every
Haagerup fusion relation and derive both traces in
\eqref{eq:defect-E8-twined}.

\par\smallskip\noindent{\footnotesize\normalfont
\hyperref[sol:ex:defect-twined]{Solution on page~\pageref*{sol:ex:defect-twined}.}}\par
\end{exercise}

\begin{exercise}\label{ex:defect-alpha}
Compute $R\alpha$, its lowest conformal weight, and its lowest
top multiplicity. Explain why this twisted sector can contain
fractional spin without violating locality of the untwisted full
theory.

\par\smallskip\noindent{\footnotesize\normalfont
\hyperref[sol:ex:defect-alpha]{Solution on page~\pageref*{sol:ex:defect-alpha}.}}\par
\end{exercise}

\begin{exercise}\label{ex:defect-between}
Compute the number of weight-$1/3$ boundary fields between two
distinct noninvertible boundaries of the extension. Compare it
with the self-boundary count nine.

\par\smallskip\noindent{\footnotesize\normalfont
\hyperref[sol:ex:defect-between]{Solution on page~\pageref*{sol:ex:defect-between}.}}\par
\end{exercise}

\begin{exercise}\label{ex:defect-categories}
Explain what is counted by 6, 12, and 36 in this chapter, and
identify which category is braided. Why would substituting local
extension modules for all extension modules erase the Haagerup
boundary category?

\par\smallskip\noindent{\footnotesize\normalfont
\hyperref[sol:ex:defect-categories]{Solution on page~\pageref*{sol:ex:defect-categories}.}}\par
\end{exercise}
